% Portable theorem guard: define only what the host preamble is missing, so this
% file drops into any of our manuscripts and re-running it is a no-op.
\makeatletter
\@ifundefined{theorem}{\newtheorem{theorem}{Theorem}}{}
\@ifundefined{lemma}{\newtheorem{lemma}[theorem]{Lemma}}{}
\@ifundefined{proposition}{\newtheorem{proposition}[theorem]{Proposition}}{}
\@ifundefined{corollary}{\newtheorem{corollary}[theorem]{Corollary}}{}
\@ifundefined{definition}{\newtheorem{definition}[theorem]{Definition}}{}
\@ifundefined{remark}{\newtheorem{remark}[theorem]{Remark}}{}
\makeatother
% BT1172 incidence-45 theorem.

\begin{theorem}[Labeled incidence model for the $45$ relation sector]
\label{thm:boolean-tritangent-incidence45}
After fixing the lexicographic bijection between the $15$ nonzero Boolean masks
and the $15$ unordered pairs of a six-set, the three-layer Boolean relation
sector carries the same share-a-line incidence graph as the cubic-surface
$45$-object model.  The graph has parameters
\[
  (v,k,\lambda,\mu)=(45,12,3,3)
\]
with $270$ edges.  Thus the Boolean relation sector is incidence-isomorphic to
the labeled three-layer cubic-surface model.  This is a labeled incidence
isomorphism; intrinsic $S_6$ naturality of the chosen mask dictionary is not
claimed here.
\end{theorem}
